\subsection{Pattern matching}
\label{sec:pmatch}

\paragraph{Presentation}

Searching for a pattern within a text is a common operation in various fields of applications, for instance for natural language processing, but also in genomics with a reduced alphabet containing the four DNA nucleotides.
We consider three state-of-the-art pattern matching algorithms: Aho-Corasick (AC), Knuth-Morris-Pratt (KMP) and Boyer-Moore (BM).
The Aho-Corasick algorithm~\cite{Aho1975} allows to search for multiple patterns while scanning the text only once, by representing the patterns in a trie and adding failure links to create an automaton.
The Knuth-Morris-Pratt algorithm~\cite{Knuth1977} uses a pre-computed longest-prefix-suffix array to avoid repeating the same comparisons between the pattern and the text.
The Boyer-Moore algorithm~\cite{Boyer1977} has a similar approach and skips many positions compared to a naive searching operation. It uses some pre-computed lookup tables to apply a bad-character rule and a good-suffix rule.

Although these algorithms only perform exact matching, they can be adapted to allow for single-letter wildcards to appear in the pattern.
We break the input pattern to obtain all slices without any wildcards, and search for the hits of each slice.
The resulting positions, offset by the starting index of the slice within the initial pattern, are then reported in a table of counters.
When a counter reaches the total amount of slices, this means there is a full match of the initial pattern at this position.

We implement the three algorithms with the BPL library.
The pre-processing (i.e. building the automaton or the lookup tables) is performed by the main processor.
The reference text is split over all processing units, with some overlapping between the fragments.
Each task receives a reference fragment, the pre-computed data, and scans its fragment to apply the algorithm and report the hit positions.

We also write an equivalent implementation that uses the main processor only (denoted Std-Multicore), and another that relies on the UPMEM's SDK (denoted SDK-PiM).
 
\paragraph{Limitations encountered}
None: the code can be fully architecture-agnostic.

\paragraph{Benchmarks}

To compare the performances of a given algorithm depending on the target architecture and implementation, we sample a reference genome and generate datasets of patterns of various characteristics.
We set the length $L$ of the patterns to 20 or 50 nucleotides.
In a dataset, we aggregate a total $N$ of 1 or 100 patterns.
Also, we introduce a number $W$ of single-letter wildcards in the pattern, either 0 or 5.
We consider two reference genomes: the human genome GRCh38~\cite{HG38} and the human gut microbiome genome UHGG~\cite{UHGG}.

\paragraph{Results} We present the execution time performances of all datasets in \tablename~\ref{tab:pmatch_results}.
For the PiM architecture, the first pattern of a dataset is slower than the following because the reference fragment is sent only at the first task execution.
The datasets of 100 patterns are thus interesting to measure the full speed the processing can achieve once the reference is set up on all processing units.
We see that the BPL-PiM version has then similar performances to the SDK-based implementation.
The three pattern matching we consider have a scanning procedure that is fairly simple, the ingenuity lying in the pre-computed data and lookup tables.
In this use case where there is little room for advanced instructions and manual optimizations on the DPUs, the BPL exhibits negligible overhead and allows to write easily a standard-like program that performs well.

\begin{table}
    \centering
    \caption{Results of the pattern matching benchmarks}
    \label{tab:pmatch_results}
    \begin{tabular}{lllllrrr}
    \toprule
        Reference & $L$ & $N$ & $W$ & Architecture & AC (s) & KMP (s) & BM (s) \\

\midrule
\multirow{4}{*}{GRCh38} & \multirow{4}{*}{20} & \multirow{4}{*}{1} & \multirow{4}{*}{0} & Std-Multicore & 1.485 & 0.872 & 0.361 \\
&&&& BPL-Multicore & 1.399 & 0.765 & 0.363 \\
\arrayrulecolor{lightgray}\cmidrule{5-8}\arrayrulecolor{black}
&&&& BPL-PiM & 0.174 & 0.255 & 0.132 \\
&&&& SDK-PiM & 0.226 & 0.258 & 0.149 \\
\midrule
\multirow{4}{*}{GRCh38} & \multirow{4}{*}{50} & \multirow{4}{*}{1} & \multirow{4}{*}{0} & Std-Multicore & 1.458 & 0.811 & 0.277 \\
&&&& BPL-Multicore & 1.348 & 0.71 & 0.276 \\
\arrayrulecolor{lightgray}\cmidrule{5-8}\arrayrulecolor{black}
&&&& BPL-PiM & 0.178 & 0.22 & 0.14 \\
&&&& SDK-PiM & 0.225 & 0.227 & 0.195 \\
\midrule
\multirow{4}{*}{GRCh38} & \multirow{4}{*}{20} & \multirow{4}{*}{100} & \multirow{4}{*}{0} & Std-Multicore & 146.954 & 81.617 & 39.711 \\
&&&& BPL-Multicore & 137.9 & 76.47 & 41.343 \\
\arrayrulecolor{lightgray}\cmidrule{5-8}\arrayrulecolor{black}
&&&& BPL-PiM & 11.152 & 18.752 & 9.143 \\
&&&& SDK-PiM & 10.446 & 17.875 & 8.998 \\
\midrule
\multirow{4}{*}{GRCh38} & \multirow{4}{*}{50} & \multirow{4}{*}{100} & \multirow{4}{*}{0} & Std-Multicore & 147.255 & 82.016 & 34.009 \\
&&&& BPL-Multicore & 139.468 & 75.447 & 37.006 \\
\arrayrulecolor{lightgray}\cmidrule{5-8}\arrayrulecolor{black}
&&&& BPL-PiM & 10.904 & 18.092 & 8.718 \\
&&&& SDK-PiM & 10.344 & 17.24 & 8.333 \\
\midrule
\multirow{4}{*}{GRCh38} & \multirow{4}{*}{50} & \multirow{4}{*}{1} & \multirow{4}{*}{5} & Std-Multicore & 2.507 & 5.314 & 3.852 \\
&&&& BPL-Multicore & 2.455 & 4.778 & 3.79 \\
\arrayrulecolor{lightgray}\cmidrule{5-8}\arrayrulecolor{black}
&&&& BPL-PiM & 1.651 & 2.304 & 1.638 \\
&&&& SDK-PiM & 1.305 & 1.986 & 1.353 \\
\midrule
\multirow{4}{*}{GRCh38} & \multirow{4}{*}{50} & \multirow{4}{*}{100} & \multirow{4}{*}{5} & Std-Multicore & 208.504 & 510.858 & 343.22 \\
&&&& BPL-Multicore & 199.283 & 456.502 & 343.321 \\
\arrayrulecolor{lightgray}\cmidrule{5-8}\arrayrulecolor{black}
&&&& BPL-PiM & 83.024 & 153.039 & 95.278 \\
&&&& SDK-PiM & 79.882 & 147.772 & 91.645 \\
\midrule
\multirow{4}{*}{UHGG} & \multirow{4}{*}{20} & \multirow{4}{*}{100} & \multirow{4}{*}{0} & Std-Multicore & 551.344 & 298.103 & 142.182 \\
&&&& BPL-Multicore & 512.282 & 267.182 & 149.261 \\
\arrayrulecolor{lightgray}\cmidrule{5-8}\arrayrulecolor{black}
&&&& BPL-PiM & 35.867 & 57.836 & 20.858 \\
&&&& SDK-PiM & 33.51 & 55.264 & 20.301 \\

\bottomrule
    \end{tabular}
\end{table}

\paragraph{Source code} The implementations are publicly available at \url{https://gitlab.inria.fr/pim/org.pim.pmatch-bpl} (GNU AGPLv3).